\section*{Chapter \ref{ch:iv:brainteasers-and-logic} --- Brainteasers and Logic}

\begin{solution}{iq:iv:brainteasers-and-logic:1}
Blue. Each move changes the number of blue tokens by 0 (two reds out, a red back; two blues out, a red back:
$-2$) or leaves it (one of each out, a blue back): the parity of the blue count never changes. It starts odd
(13), so the last token, when the bag holds one, is blue. With 12 blues it would be red.
\iqlookfor{the parity invariant, found by writing out the three moves.}
\end{solution}

\begin{solution}{iq:iv:brainteasers-and-logic:2}
The losing positions for the player to move are the multiples of 5: from a multiple of 5 every move leaves a
non-multiple, and from a non-multiple a move to a multiple exists. With 30 tokens, go second. With 31, take one
and then always answer the opponent's $k$ with $5 - k$.
\iqlookfor{\omterm{def:iv:brainteasers-and-logic:backward}{backward induction} to a periodic pattern and the pairing strategy.}
\end{solution}

\begin{solution}{iq:iv:brainteasers-and-logic:3}
There are 366 possible birthdays (29 February included) and 400 people; by the pigeonhole principle two share
one. The argument needs no assumption about how birthdays are distributed.
\iqlookfor{naming the boxes and counting.}
\end{solution}

\begin{solution}{iq:iv:brainteasers-and-logic:4}
Summing each trader's number of counterparties counts every pair twice, so the sum is even. A sum of integers is
even only if the number of odd terms is even. (A check on random graphs in the chapter's code agrees.)
\iqlookfor{double counting and a parity argument.}
\end{solution}

\begin{solution}{iq:iv:brainteasers-and-logic:5}
Start both at 0. At 3 minutes the small timer runs out: turn it over. At 5 minutes the large one runs out; the
small one has run for 2 minutes since it was turned, so turn it over again: its 2 minutes of sand run back, ending
at 7 minutes.
\iqlookfor{tracking the sand in both halves, not only the elapsed time.}
\end{solution}

\begin{solution}{iq:iv:brainteasers-and-logic:6}
Number the colours 0, 1, 2. The last person announces the sum of the hats in front modulo 3 (a colour, so a
legal answer). Each following person knows that sum, subtracts the colours announced by those behind them after
the first (which are correct) and the colours they see in front, and obtains their own colour modulo 3. Everyone
but the last is right: at least 99. The same rule works for any number of colours, which is the family's idea:
one announcement can carry one number modulo $k$.
\iqlookfor{the sum-modulo-$k$ invariant, and the generalisation.}
\end{solution}

\begin{solution}{iq:iv:brainteasers-and-logic:7}
Take $i$ coins from bag $i$, 55 coins in all, and weigh them. If every coin were genuine the scale would read
55 grams; the reading is short by $0.1\,i$ grams when bag $i$ is light, so the shortfall divided by 0.1 is the bag's
number. The design gives each hypothesis a distinct, readable outcome.
\iqlookfor{encoding the hypothesis in the quantity taken.}
\end{solution}

\begin{solution}{iq:iv:brainteasers-and-logic:8}
One partner keeps everything. With two, the senior's vote is half, so it passes: $(100, 0)$. With three, the
proposal needs one vote besides the proposer's: the partner who would get nothing in the two-partner case (the
most junior) is bought with 1: $(99, 0, 1)$. With four, two votes of four pass, so again one vote is needed
besides the proposer's. If the proposal failed, the other three would split $(99, 0, 1)$: the second partner would
get 99, the third 0 and the fourth 1. The cheapest vote is the third partner's, bought with 1: the proposal is
$(99, 0, 1, 0)$. An exhaustive
\omterm{def:iv:brainteasers-and-logic:backward}{backward induction} in the chapter's code gives the same, and $(98, 0, 1, 0, 1)$ for five.
\iqlookfor{\omterm{def:iv:brainteasers-and-logic:backward}{backward induction} from the smallest case and buying the cheapest votes.}
\end{solution}

\begin{solution}{iq:iv:brainteasers-and-logic:9}
The game is finite and cannot end in a draw, so one player has a winning strategy. Suppose it is the second.
Let the first player bite only the bottom-right square. Whatever the second player answers is a position the
first player could have produced in one bite from the full bar, because any bite also removes the bottom-right
square. So the first player could have made that bite first and then used the second player's winning strategy:
a contradiction. The first player wins; an exhaustive search on the four-by-five bar confirms it, and finds a
move, but the argument never needs one.
\iqlookfor{the \omterm{def:iv:brainteasers-and-logic:stealing}{strategy-stealing argument}, including why the stolen position is reachable.}
\end{solution}

\begin{solution}{iq:iv:brainteasers-and-logic:10}
The three red-badged traders step forward together at the third round. If there were one red badge, its wearer
would see none and know at once. If there were two, each would see one; when nobody steps forward at round one,
each learns that the other saw a red badge, which must be their own, so both step forward at round two. With
three, each red wearer sees two and waits; when nobody steps forward at round two, each concludes that there are
three. The sentence added common knowledge: before it, everyone knew there was a red badge, but not that everyone
knew that everyone knew it, which is what the induction needs. A model-checking search over all consistent
worlds gives round 3 and the three red wearers.
\iqlookfor{induction on the number of red badges and the idea of common knowledge.}
\end{solution}

\begin{solution}{iq:iv:brainteasers-and-logic:11}
Rule: if the two hats you see have the same colour, guess the other colour; otherwise pass. The team loses only
when all three hats match (two of the eight configurations), where everyone guesses wrong; in the other six exactly
one player sees two equal hats and guesses right: $\tfrac34$. No strategy does better: for each player and each view,
a guess is right in one of the two configurations sharing that view and wrong in the other, so the total number of
right guesses equals the total of wrong ones. A win needs at least one right guess and a loss can absorb up to three
wrong ones, so wins are at most three times losses: at most $\tfrac34$. With seven players a Hamming code wins
$\tfrac78$ of the time; both are checked exhaustively in the chapter's code.
\iqlookfor{concentrating the wrong guesses on few configurations, and the counting bound.}
\end{solution}

\begin{solution}{iq:iv:brainteasers-and-logic:12}
$1 + (a + b + ab) = (1 + a)(1 + b)$, so the product $\prod(1 + x)$ over the board is invariant. It starts at
$2 \times 3 \times \dots \times 11 = 11!$, so the last number is $11! - 1 = 39\,916\,799$, whatever the order.
\iqlookfor{spotting the factorisation that turns the move into a product.}
\end{solution}

\begin{solution}{iq:iv:brainteasers-and-logic:13}
Computing losing positions upwards with moves $\{1, 3, 4\}$ gives 0, 2, 7, 9, 14, 16, 21, 23, \dots: the
positions congruent to 0 or 2 modulo 7. Twenty is winning: take 4, leaving 16. The pattern is periodic, but its
period (7) is not $1 + $ the largest move, which is why it must be computed rather than guessed.
\iqlookfor{a clean backward-induction table and suspicion of patterns guessed from the classic game.}
\end{solution}

\begin{solution}{iq:iv:brainteasers-and-logic:14}
The prefix sums $P_0 = 0, P_1, \dots, P_{10}$ are eleven numbers in ten residue classes modulo 10, so two are
congruent, and the days strictly between them sum to a multiple of 10. Ten is the fewest: nine days of P\&L equal
to 1 have runs summing to 1 to 9 only. A check on 500 random sequences of ten confirms the statement.
\iqlookfor{the pigeonhole on prefix sums and a tightness example.}
\end{solution}
